
\section*{6. The \(p\)-Lie Algebra of an \(S\)-Group Scheme}
\phantomsection
\label{sga3:VIIA:6}
\addcontentsline{toc}{section}{6. The p-Lie Algebra of an S-Group Scheme}

Let \(S\) be a scheme of characteristic \(p>0\). In \SGARefLink{sga3:VIIA:5.5}{5.5} we
associated with every quasi-coherent
\(\mathcal O_S\)-\(p\)-Lie algebra \(\mathcal L\) the group
\(S\)-functor
\[
  \mathfrak G_p(\mathcal L)
  =
  \operatorname{Spec}^{*}\mathcal U_p(\mathcal L).
\]
We shall now see that, for every \(S\)-group scheme \(G\), the
\(\mathcal O_S\)-Lie algebra \(\operatorname{Lie}(G/S)\) defined
in II, \SGARefLink{sga3:ii:subsection:4-11}{4.11} is naturally endowed with the structure of an
\(\mathcal O_S\)-\(p\)-Lie algebra.

\medskip
\noindent\textbf{6.1.}
\phantomsection
\label{sga3:VIIA:6.1}

First identify \(\operatorname{Lie}(G/S)(S)\) and
\(\operatorname{Lie}(\operatorname{Aut}G/S)(S)\), respectively,
with Lie subalgebras of \(\mathrm U(G)\) and
\(\operatorname{Dif}_{G/S}\) by means of the injections
\(\alpha\) and \(\beta\) of \SGARefLink{sga3:VIIB:2.5:proposition}{2.5}. Thus
\(\operatorname{Lie}(\operatorname{Aut}G/S)(S)\) is identified
with the \(\Gamma(S,\mathcal O_S)\)-Lie algebra of
\(S\)-derivations of \(\mathcal O_G\). By \SGARefLink{sga3:VIIB:5.2:theorem}{5.2}, the latter is a
restricted \(p\)-Lie subalgebra of
\(\operatorname{Dif}_{G/S}\).

On the other hand, the image of
\[
  L=\operatorname{Lie}(G/S)(S)
\]
under the injective algebra homomorphism
\[
  \ell:\mathrm U(G)\longrightarrow\operatorname{Dif}_{G/S},
  \qquad
  d\longmapsto {}_Gd,
\]
consists of the left-invariant derivations (cf.\ \SGARefLink{sga3:VIIA:2.2}{2.2}, \SGARefLink{sga3:VIIA:2.4}{2.4}, and \SGARefLink{sga3:VIIB:2.5:proposition}{2.5}). If \(x\in L\), then
\(\ell(x)^p=\ell(x^p)\) by the cited passages. Since
\(\ell(x)^p\) is again a derivation, \(x^p\) belongs to
\(\operatorname{Lie}(G/S)(S)\). Hence
\begin{center}
  \textit{\(\operatorname{Lie}(G/S)(S)\) is a restricted
  \(p\)-Lie subalgebra of the infinitesimal algebra
  \(\mathrm U(G)\).}
\end{center}

\medskip
\noindent\textbf{6.1.1.}
\phantomsection
\label{sga3:VIIA:6.1.1}

Let \(\varphi:G\to H\) be a homomorphism of \(S\)-group schemes.
The homomorphisms
\(\operatorname{Lie}(\varphi/S)(S)\) and \(\mathrm U(\varphi)\)
are compatible with the identifications of
\(\operatorname{Lie}(G/S)(S)\) and
\(\operatorname{Lie}(H/S)(S)\) with restricted \(p\)-Lie
subalgebras of \(\mathrm U(G)\) and \(\mathrm U(H)\). Since
\(\mathrm U(\varphi)\) is an algebra homomorphism,
\(\operatorname{Lie}(\varphi/S)(S)\) is a homomorphism of
restricted \(p\)-Lie algebras.

Likewise, if \(s:T\to S\) is a base change, the map from
\(\operatorname{Lie}(G/S)(S)\) to
\(\operatorname{Lie}(G/S)(T)\) induced by \(s\) is a
homomorphism of restricted \(p\)-Lie algebras. Equivalently, the
functor \(\operatorname{Lie}(G/S)\) is endowed with a structure
of \(\mathcal O_S\)-\(p\)-Lie algebra. In particular, as \(T\)
ranges over the open subsets of \(S\), the sheaf
\(\operatorname{Lie}(G/S)\) is endowed with a structure of
\(\mathcal O_S\)-\(p\)-Lie algebra.

\medskip
\noindent\textbf{6.2.}
\phantomsection
\label{sga3:VIIA:6.2}

Following an idea of Demazure, we shall generalize the preceding
construction to certain group \(S\)-functors which need not be
representable. We first give another definition of the symbolic
\(p\)-th power in the Lie algebra of an \(S\)-group scheme \(G\).

Let \(D\) be a derivation of \(G\) at the
origin.\footnote{The editors expanded the argument in the
original text in what follows.} By \SGARefLink{sga3:VIIA:1.2.1}{1.2.1}, \(D\) is the composite
of the \(S\)-derivation
\[
  \delta=(\tau,\partial_t)
\]
of the zero section \(\tau:S\to I_S\), and a morphism
\(x:I_S\to G\) such that \(x\circ\tau=\varepsilon\), that is,
\(x\in\operatorname{Lie}(G/S)(S)\). By the definition in \SGARefLink{sga3:VIIA:2.1}{2.1},
\(D^p\) is the following composite deviation:
\begin{center}
  \phantomsection
  \label{sga3:VIIA:6.2-demazure-pfold-composite}
  \resizebox{0.98\linewidth}{!}{%
  \(\displaystyle
  S \simeq
  \underbrace{S\times S\times\cdots\times S}_{p}
    \xrightarrow{\delta\times\cdots\times\delta}
  I_S\times\cdots\times I_S
    \xrightarrow{x\times\cdots\times x}
  G\times\cdots\times G
    \xrightarrow{m^{(p)}} G
  \)%
}

\end{center}
Here \(m^{(p)}\) is induced by the multiplication
\(m:G\times G\to G\). Since
\(I_S\times\cdots\times I_S\) is affine over \(S\), with affine
algebra
\[
  \mathcal B
  =
  \mathcal O_S[d_1,\ldots,d_p]/
  (d_1^2,\ldots,d_p^2),
\]
the deviation \(\delta\times\cdots\times\delta\) is defined by
the homomorphism of \(\mathcal O_S\)-modules
\[
  \phi:\mathcal B\longrightarrow\mathcal O_S
\]
which sends the monomial \(d_1d_2\cdots d_p\) to \(1\), and all
the other monomials \(d_{i_1}\cdots d_{i_r}\), for
\(0\le r<p\), to \(0\).

Let \(\operatorname{pr}_i\) denote the projection of
\(I_S^p\) onto its \(i\)-th factor, and let \(x_i\) be the image
of \(x\) in \(G(I_S^p)\) under \(G(\operatorname{pr}_i)\). Then
\[
  m^{(p)}\circ(x\times\cdots\times x)
\]
is precisely the product \(x_1x_2\cdots x_p\). Consequently,
\(D^p\) is also the composite deviation
\begin{center}
  \phantomsection
  \label{sga3:VIIA:6.2-demazure-product-composite}
  \(\displaystyle
S \xrightarrow{\delta\times\cdots\times\delta}
I_S\times\cdots\times I_S
\xrightarrow{x_1x_2\cdots x_p}
G .
\)

\end{center}

This description gives another proof that \(D^p\) is a
derivation of \(G\) at the origin. Since \(G\) is a very good
group (II, \SGARefLink{sga3:ii:subsection:4-11}{4.11}), the images \(G(\operatorname{pr}_1)(x)\) and
\(G(\operatorname{pr}_2)(x)\) of \(x\) in
\(G(I_S\times I_S)\) commute. It follows that the elements
\(x_i\in G(I_S^p)\) commute pairwise. Hence, for every
permutation \(\gamma\) of the factors of \(I_S^p\),
\[
  (x_1\cdots x_p)\circ\gamma=x_1\cdots x_p.
\]
Thus \(x_1\cdots x_p\) factors through the canonical projection
from \(I_S^p\) to the symmetric product \(\Sigma^p I_S\)
(cf.\ \SGARefLink{sga3:VIIA:4.2}{4.2}).

The affine algebra of \(\Sigma^p I_S\) is the subalgebra
\(\mathcal A\) of \(\mathcal B\) having as an
\(\mathcal O_S\)-basis the elementary symmetric functions
\[
  1=\sigma_0,\sigma_1,\ldots,\sigma_p
\]
in \(d_1,\ldots,d_p\). Let
\(\kappa:\mathcal A\hookrightarrow\mathcal B\) be the
inclusion, and let
\[
  \pi:\mathcal A\longrightarrow\mathcal O_S[t]/(t^2)
\]
be the homomorphism of \(\mathcal O_S\)-algebras which kills
\(\sigma_1,\ldots,\sigma_{p-1}\) and sends \(\sigma_p\) to
\(t\). Then
\[
  \phi\circ\kappa=\partial_t\circ\pi,
\]
where \(\partial_t:\mathcal O_S[t]/(t^2)\to\mathcal O_S\)
kills \(1\) and sends \(t\) to \(1\). If
\[
  i:I_S\hookrightarrow\Sigma^p I_S
\]
denotes the closed immersion defined by \(\pi\), we obtain the
commutative diagram
\begin{center}
  \phantomsection
  \label{sga3:VIIA:6.2-symbolic-power-commutative-diagram}
  \begin{tikzcd}[row sep=3.5em,column sep=4.3em]
S
  \arrow[r,swap,"\delta\times\cdots\times\delta"]
  \arrow[d,swap,"\delta"]
  \arrow[rr,bend left=23,"D^p"]
  & I_S^p \arrow[r,swap,"x_1\cdots x_p"] \arrow[d,"\mathrm{can.}"]
  & G \arrow[d,equal] \\
I_S
  \arrow[r,"i"]
  \arrow[rr,bend right=23,swap,"y"]
  & \Sigma^p I_S \arrow[r]
  & G
\end{tikzcd}

\end{center}
which shows that \(D^p\) has the form \(y\circ\delta\), and is
therefore indeed a derivation of \(G\) at the origin.

\medskip
\noindent\textbf{6.3.}
\phantomsection
\label{sga3:VIIA:6.3}

Let \(\mathfrak S_p\) be the symmetric group of degree \(p\),
and let \(I_S^p\times\mathfrak S_p\) denote the disjoint union
of a family of copies of \(I_S^p\) indexed by
\(\mathfrak S_p\). Let
\[
  \pi:I_S^p\times\mathfrak S_p\longrightarrow I_S^p
\]
be the canonical projection, and let
\[
  \mu:I_S^p\times\mathfrak S_p\longrightarrow I_S^p
\]
be the morphism defining the action of \(\mathfrak S_p\) on
\(I_S^p\): if \(\rho\in\mathfrak S_p\), the restriction of
\(\mu\) to \(I_S^p\times\rho\) has
\(\operatorname{pr}_{\rho(j)}\) as its \(j\)-th component.

\noindent\textbf{Definition.}
A functor
\[
  X:(\operatorname{Sch}/S)^\circ\longrightarrow(\operatorname{Ens})
\]
satisfies \emph{condition (F)} if:
\begin{enumerate}
  \item[\textit{a)}] \(X\) takes finite disjoint unions to direct
  products;
  \item[\textit{b)}] for every \(S\)-scheme \(T\), the following
  sequence is exact:
\end{enumerate}
\begin{center}
  \phantomsection
  \label{sga3:VIIA:6.3-conditionF-equalizer}
  \resizebox{0.98\linewidth}{!}{%
\begin{tikzcd}[ampersand replacement=\&,column sep=4.4em]
X(T\times\Sigma^p I_S)
  \arrow[r]
  \& X(T\times I_S^p)
    \arrow[r,shift left=.8ex,"X(\operatorname{id}_T\times\pi)"]
    \arrow[r,shift right=.8ex,swap,"X(\operatorname{id}_T\times\mu)"]
  \& X(T\times I_S^p\times\mathfrak S_p)\, .
\end{tikzcd}%
}

\end{center}

Every \(S\)-scheme satisfies (F). If \(\mathcal F\) is an
\(\mathcal O_S\)-module, then \(\mathcal W(\mathcal F)\)
satisfies (F). Every projective limit of functors satisfying
(F) again satisfies (F). Finally, if \(Y\) satisfies (F) and
\(X\) is an arbitrary \(S\)-functor, then
\(\operatorname{Hom}_S(X,Y)\) satisfies (F).

Let \(X\) be a very good group (II, \SGARefLink{sga3:ii:definition:4-10}{4.10}) satisfying condition
(F). Let \(x:I_S\to X\) be a morphism extending the unit section
of \(X\), and use the notation of \SGARefLink{sga3:VIIA:6.2}{6.2}. As above,
\[
  x_1\cdots x_p:I_S^p\longrightarrow X
\]
factors through \(\Sigma^p I_S\):
\begin{center}
  \phantomsection
  \label{sga3:VIIA:6.3-symbolic-power-factorization}
  \begin{tikzcd}[row sep=3.5em,column sep=2.1em]
I_S^p \arrow[rr,"x_1\cdots x_p"] \arrow[dr,swap,"\mathrm{can.}"]
  && X \\
& \Sigma^p I_S \arrow[ur,swap,"\Sigma^p(x)"]
&
\end{tikzcd}

\end{center}
and therefore defines, by composition, a morphism
\begin{center}
  \phantomsection
  \label{sga3:VIIA:6.3-symbolic-pth-power-composite}
  \(\displaystyle
x^{(p)}:\quad
I_S \xrightarrow{i} \Sigma^p I_S
  \xrightarrow{\Sigma^p(x)} X
\)

\end{center}
which we call the \emph{symbolic \(p\)-th power} of \(x\).

The endomorphism
\[
  x\longmapsto x^{(p)}
\]
of \(\operatorname{Lie}(X/S)(S)\) is plainly compatible with
base change and functorial in \(X\). It would be interesting to
know for which \(X\) this endomorphism makes
\(\operatorname{Lie}(X/S)(S)\) into a restricted \(p\)-Lie
algebra.

\medskip
\noindent\textbf{6.4.}
\phantomsection
\label{sga3:VIIA:6.4}

The preceding definition of the symbolic \(p\)-th power is
particularly well suited to computation. Here are some
examples.

\medskip
\noindent\textbf{6.4.1.}
\phantomsection
\label{sga3:VIIA:6.4.1}

Let \(M\) be an ``abstract'' abelian group, and let \(D_S(M)\)
be the diagonalizable \(S\)-group of type \(M\) (I, \SGARefLink{sga3:i:paragraph:4-4-2}{4.4.2}). For
every \(S\)-scheme \(T\),
\[
  D_S(M)(T)
  =
  \operatorname{Hom}_{(\mathrm{Ab})}
  \bigl(M,\mathcal O(T)^\times\bigr).
\]
Let \(x\in\operatorname{Lie}(D_S(M)/S)(S)\), that is, an
abelian-group homomorphism
\[
  x:M\longrightarrow
  \Gamma(S,\mathcal O_S+d\mathcal O_S)^\times
\]
of the form
\[
  m\longmapsto1+d\xi(m),
  \qquad
  \xi\in
  \operatorname{Hom}_{(\mathrm{Ab})}(M,\mathcal O(S)).
\]
With the notation of \SGARefLink{sga3:VIIA:6.2}{6.2} and \SGARefLink{sga3:VIIA:6.3}{6.3}, the product
\(x_1\cdots x_p\) associates with \(m\in M\) the expression
\[
  (1+d_1\xi(m))\cdots(1+d_p\xi(m)),
\]
that is,
\[
  1+\sigma_1\xi(m)+\sigma_2\xi(m)^2+\cdots+
  \sigma_p\xi(m)^p.
\]
This expression belongs to \(\mathcal O(\Sigma^p I_S)\).
Projecting it to \(\mathcal O(S)[d]/(d^2)\), killing
\(\sigma_1,\ldots,\sigma_{p-1}\), and sending \(\sigma_p\) to
\(d\), we see that \(x^{(p)}\) is the homomorphism
\[
  m\longmapsto1+d\xi(m)^p.
\]
Thus, under the identification
\[
  \operatorname{Lie}(D_S(M)/S)(S)
  \simeq
  \operatorname{Hom}_{(\mathrm{Ab})}(M,\mathcal O(S))
\]
of II, \SGARefLink{sga3:ii:proposition:5-1}{5.1}, the symbolic \(p\)-th power sends \(\xi\) to
\[
  \xi^{(p)}:m\longmapsto\xi(m)^p.
\]

\medskip
\noindent\textbf{6.4.2.}
\phantomsection
\label{sga3:VIIA:6.4.2}

Let \(\mathcal F\) be an \(\mathcal O_S\)-module, and let
\(G=\mathcal W(\mathcal F)\) be the group \(S\)-functor in
abelian groups (cf.\ I, \SGARefLink{sga3:i:subsection:4-6}{4.6}). Let
\[
  y\in\mathcal W(\mathcal F)(S)=\Gamma(S,\mathcal F),
\]
and let \(y'\) be its image in
\(\mathcal W(\mathcal F)(I_S)\) under
\(\mathcal W(\mathcal F)(I_S\to S)\).

By II, \SGARefLink{sga3:ii:examples:4-4-2}{4.4.2} and \SGARefLink{sga3:ii:scholium:4-5-1}{4.5.1}, the map
\[
  y\longmapsto dy'
\]
is an isomorphism of \(\mathcal O(S)\)-modules from
\(\mathcal W(\mathcal F)(S)\) onto
\(\operatorname{Lie}(\mathcal W(\mathcal F)/S)(S)\). Put
\(x=dy'\). Then the element \(x_i\) of \SGARefLink{sga3:VIIA:6.2}{6.2} is
\(d_i y''\), where \(y''\) denotes the canonical image of
\(y'\) in
\(\mathcal W(\mathcal F)(I_S^p)\). Hence
\[
  x_1+\cdots+x_p
  =
  (d_1+\cdots+d_p)y''
  =
  \sigma_1y'',
\]
and this belongs to
\(\mathcal W(\mathcal F)(\Sigma^p I_S)\). Since the
homomorphism
\[
  \mathcal O(\Sigma^p I_S)\longrightarrow\mathcal O(I_S)
\]
defining the morphism \(i\) kills \(\sigma_1\), we have
\(x^{(p)}=0\). Therefore
\begin{center}
  \textit{For every \(\mathcal O_S\)-module \(\mathcal F\), the
  operation \(x\mapsto x^{(p)}\) on
  \(\operatorname{Lie}\mathcal W(\mathcal F)\) is zero.}
\end{center}

\medskip
\noindent\textbf{6.4.3.}
\phantomsection
\label{sga3:VIIA:6.4.3}

Let \(X\) be an \(S\)-scheme, let
\(G=\operatorname{Aut}_S X\), and let \(D\) be an
\(S\)-derivation of the structure sheaf \(\mathcal O_X\). By
\SGARefLink{sga3:VIIA:1.5}{1.5}, \(D\) may be identified with an \(I_S\)-automorphism \(x\)
of \(X_{I_S}\) inducing the identity on \(X\). It can be
described as follows. If \(f\) is a section of
\(\mathcal O_X[d]/(d^2)\) of the form \(a+bd\), put
\[
  D_{I_S}f=Da+(Db)d.
\]
Thus \(D_{I_S}\) is obtained from \(D\) by the base change
\(I_S\to S\), and the automorphism of \(X_{I_S}\) is associated
with the endomorphism
\[
  f\longmapsto f+(D_{I_S}f)d
  =
  a+(D(a)+b)d
\]
of \(\mathcal O_X[d]/(d^2)\).

Similarly, let \(D_{I_S^p}\) be the differential operator on
\(X_{I_S^p}\) obtained from \(D\) by the base change
\(I_S^p\to S\). With the notation of \SGARefLink{sga3:VIIA:6.2}{6.2}, the automorphism
\(x_i\) of \(X_{I_S^p}\) is associated with the endomorphism
\[
  f\longmapsto f+(D_{I_S^p}f)d_i
\]
of
\[
  \mathcal O_X[d_1,\ldots,d_p]/(d_1^2,\ldots,d_p^2).
\]
The product \(x_1\cdots x_p\) is therefore associated with
\[
  (1+d_1D_{I_S^p})(1+d_2D_{I_S^p})\cdots
  (1+d_pD_{I_S^p}),
\]
that is, with
\[
  1+\sigma_1D_{I_S^p}
  +\sigma_2D_{I_S^p}^{\,2}
  +\cdots+
  \sigma_pD_{I_S^p}^{\,p}.
\]
The coefficient of \(\sigma_p\) is \(D_{I_S^p}^{\,p}\).
Consequently, the Lie-algebra isomorphism
\[
  \operatorname{Der}_S(\mathcal O_X)
  \xrightarrow{\ \sim\ }
  \operatorname{Lie}(\operatorname{Aut}_S X),
  \qquad
  D\longmapsto x
\]
(cf.\ \SGARefLink{sga3:VIIA:1.5}{1.5}) is also an isomorphism of restricted \(p\)-Lie
algebras.

\medskip
\noindent\textbf{6.4.4.}
\phantomsection
\label{sga3:VIIA:6.4.4}

The same method shows that, for every \(\mathcal O_S\)-module
\(\mathcal F\), the Lie-algebra isomorphism
\[
  \operatorname{Lie}
  \bigl(
    \operatorname{Aut}_{\mathcal O_S\text{-mod.}}
    \mathcal W(\mathcal F)/S
  \bigr)(S)
  \xrightarrow{\ \sim\ }
  \bigl(
    \operatorname{End}_{\mathcal O_S\text{-mod.}}
    \mathcal W(\mathcal F)
  \bigr)(S)
\]
(cf.\ II, \SGARefLink{sga3:ii:proposition:4-8}{4.8}) is also an isomorphism of restricted
\(p\)-Lie algebras.

\medskip
\noindent\textbf{6.4.5.}\footnote{The editors supplied the
number \SGARefLink{sga3:VIIA:6.4.5}{6.4.5} and expanded the original text.}
\phantomsection
\label{sga3:VIIA:6.4.5}

Let \(\mathcal U\) be an \(\mathcal O_S\)-coalgebra in groups,
and let
\[
  G=\operatorname{Spec}^{*}\mathcal U
\]
be the group \(S\)-functor; suppose that \(G\) is representable.
For every \(T\to S\), \SGARefLink{sga3:VIIA:5.5}{5.5} and \SGARefLink{sga3:VIIA:6.1.1}{6.1.1} define two restricted
\(p\)-Lie algebra structures on
\[
  L(T)=\operatorname{Lie}(G)(T).
\]
There is a commutative diagram
\begin{center}
  \phantomsection
  \label{sga3:VIIA:6.4.5-group-coalgebra-comparison}
  \begin{tikzcd}[row sep=3.4em,column sep=4.3em]
L(T) \arrow[r,hook,"\tau"] \arrow[d,hook,swap,"\alpha"] \arrow[dr,"i"]
  & \Gamma(T,\mathcal U_T) \\
\mathrm U(G_T)
  & \mathrm U(L(T)) \arrow[l,swap,"\varphi"] \arrow[u,swap,"\psi"]
\end{tikzcd}

\end{center}
in which \(\mathrm U(L(T))\) is the enveloping algebra of
\(L(T)\), and \(\varphi,\psi\) are the algebra homomorphisms
induced by \(\alpha,\tau\). It follows that the two restricted
\(p\)-Lie algebra structures coincide: if
\(x\in L(T)\) is identified with its image in
\(\mathrm U(G_T)\), respectively in
\(\Gamma(T,\mathcal U_T)\), then \(x^{(p)}\) is the image under
\(\varphi\), respectively under \(\psi\), of the element
\(x^p\in\mathrm U(L(T))\).

